% ID: SEL-D0056-OB-E01
% ESTADO: cotejada. Segunda pasada visual y borrador compilado revisados.
% Fuente: Selectividad/sel_2011_matematicas_aplicadas/123-EXAMEN 1.pdf
% SHA-256: 69f8b84ed22f2b49074eafb95dbdce8a98114a4a2f02989b92f3f7add9355155
% Página: 2; opción B; ejercicio original 1.

\textbf{(2.5 puntos)} Una empresa elabora dos productos, A y B. Cada unidad de A requiere 2 horas en una máquina y 5 horas en una segunda máquina. Cada unidad de B necesita 4 horas en la primera máquina y 3 horas en la segunda máquina. Semanalmente se dispone de 100 horas en la primera máquina y de 110 horas en la segunda.

Si la empresa obtiene un beneficio de 70 euros por cada unidad de A, y de 50 euros por cada unidad de B, \textquestiondown{}qué cantidad semanal de cada producto debe producir con objeto de maximizar el beneficio total? \textquestiondown{}Cuál es ese beneficio?
